\scanpage{70}
\mspage{II.18}
% archive sheet number: 69
% lines 1–3
\deleted{pourrait généraliser substantiellement \deleted{\unsure{le}} 1.3, 1.5, \deleted{\illegible{1}}, \deleted{\illegible{1}} 1.9 \unsure{$\#$} (ss-réserve de $\frac12$ simplicité), 1.11, en y remplaçant partout ``ess. abélienne'' par ``ess.$^{\mathrm t}$ résoluble''.}

% line 4
2. \uline{Démonstration} \inserted{de 1.1} : \uline{Cas où $X$ est de dim} 1.

% lines 5–9
\deleted{Alors $X$ est un ouvert}

2.1. Notons d'abord (sans hyp. sur la dim de $X$) qu'on peut supposer $X$ \uline{affine}. \deleted{\unsure{du schéma géom. irréduct.}} En effet, si $U$ est un ouvert affine \deleted{\unsure{de}} de $X$ contenant l'image de $\unsure{\xi|\eta}$, on a \deleted{On sait que $\pi_1($} \inserted{un \unsure{hom.} de suites exactes}
% lines 10–12
\diagram{\begin{tikzcd}[ampersand replacement=\&]
1 \arrow[r,dash] \& \pi_1(\bar U,\xi) \arrow[r,dash] \arrow[d,"\alpha"] \& \pi_1(\deleted{\unsure{X}}U,\xi) \arrow[r,dash] \arrow[d,"\beta"] \& \pi_1(S',\xi) \arrow[r,dash] \arrow[d,Rightarrow] \& 1 \\
1 \arrow[r,dash] \& \pi_1(\bar X,\xi) \arrow[r,dash] \& \pi_1(X,\xi) \arrow[r,dash] \& \pi_1(S,\xi) \arrow[r,dash] \& 1 .
\end{tikzcd}}

% lines 13–17
Comme \deleted{$X$ est} $\bar X$ \deleted{\unsure{est}} \deleted{géométriquement \unsure{unibranche}} \inserted{normal}, \deleted{(\unsure{ainsi} $\bar X$ cf. EGA IV 6.15.6)}, il s'ensuit que $\beta$ est surjectif (SGA1 V \unsure{\ldots}), donc d'après \deleted{$\pi_1(X$ $\alpha$ et $\beta$ sont surjectifs.} On en conclut que si la \deleted{\unsure{même}} conclusion de 1.1. est vraie pour $(U,\xi)$, elle est vraie pour $(X,\xi)$.

% lines left margin beside 17–24
% margin note written sideways
% slanted
\margin{\deleted{2.2. On voit de même que la conclusion \unsure{ne change pas} si on \unsure{remplace} $K$ par une \unsure{ext.} finie $K'$, et $X$ par $X\otimes_K K'$, et $\xi$ par $\xi'$ \unsure{au-dessus de} $\xi$ : cela revient à \unsure{remplacer} $\pi_1(X,\xi)$ par un ss-groupe ouvert, \unsure{l'image réciproque d'un} ss-groupe ouvert de $\pi_1(S,\xi)$.}}

% lines 18–24
\deleted{2.3. \deleted{\unsure{Soit donc}} Supposons donc \deleted{\unsure{X}} $X$ affine, donc quasi-projectif, et supposons de plus $\dim X=1$. Alors $X$ est un \deleted{\unsure{normal|sous-sch}} ouvert d'une courbe normale projective $\hat X$ (EGA \unsure{II|III} \unsure{7.1|7.}). Quitte à passer à une extension \deleted{radicielle} finie de $K$, on est ramené au cas où $\hat X'$ est \uline{lisse} sur $K$, et où les points de \deleted{\unsure{X}} $\hat X-X$ sont tous rationnels sur $K$. Soit $x_i$ ($1\leq i\leq n$) ces pts.}
